\chapter{A Concrete Naming Certificate for $\ChoiceFreeLimitScheduleUp$}
\label{ch:concrete-instances-choice-free-limit-schedule-namecert}
\origin{human}

The $\ChoiceFreeLimitScheduleUp$ packet records the finite precision-to-window route used by Real-completion consumers before a $\RealUp$ seal is reached. It keeps the precision request, schedule row, selected $\StreamNameUp$ windows, $\DyadicRatCoreUp$ observations, $\RegSeqRatUp$ handoff, $\RealUp$ seal, transports, routes, provenance, and local naming row in one displayed $\BHist$ packet, so the limit-facing surface is finite and contains no choice operator.

\begin{definition}[Choice-free limit schedule carrier]
\label{def:choice-free-limit-schedule-carrier}
A $\ChoiceFreeLimitScheduleUp$ carrier is a finite $\BHist$ packet
$$
L=(E,M,W,D,R,S,H,C,P,N)
$$
with the following displayed rows:
\begin{enumerate}
\item $E$ is the precision request row naming the finite tolerance level demanded by a completion-facing consumer;
\item $M$ is the schedule row assigning the request $E$ to the finite observation route used at that tolerance;
\item $W$ is the finite family of $\StreamNameUp$ window rows inspected by the schedule;
\item $D$ records the $\DyadicRatCoreUp$ observations and radius ledgers read from the windows in $W$;
\item $R$ records the $\RegSeqRatUp$ handoff rows reached after those dyadic observations;
\item $S$ records the $\RealUp$ seal row that receives the handoff;
\item $H$ records componentwise $\hsame$ transports for the request, schedule, windows, dyadic observations, regular-sequence handoff, seal, provenance, and local name;
\item $C$ records the displayed $\Cont$ routes from $E$ to $M$, from $M$ to $W$, from $W$ to $D$, from $D$ to $R$, and from $R$ to $S$;
\item $P$ is the $\Pkg$ provenance over $\DyadicRatCoreUp$, $\StreamNameUp$, $\RegSeqRatUp$, $\RealUp$, $\BHist$, $\hsame$, $\Cont$, and $\NameCert$;
\item $N$ is the local $\NameCert_{\ChoiceFreeLimitScheduleUp}$ row.
\end{enumerate}
The carrier is exactly this finite packet. It has no coordinate for quotient streams, host equality of completed reals, countable choice, a hidden limit selector, or an ambient completeness theorem.
\end{definition}

\begin{theorem}[Choice-free limit schedule NameCert obligations]
\label{thm:choice-free-limit-schedule-namecert-obligations}
The local $\NameCert_{\ChoiceFreeLimitScheduleUp}$ surface for \autoref{def:choice-free-limit-schedule-carrier} consists of five finite obligations:
\begin{enumerate}
\item carrier exposure, requiring every accepted object to display the packet $L=(E,M,W,D,R,S,H,C,P,N)$;
\item request-to-schedule exactness, requiring the precision request $E$ to reach the schedule row $M$ only through the displayed route in $C$;
\item window and dyadic readback coverage, requiring $M$ to select only the finite $\StreamNameUp$ windows in $W$ and requiring every exported $\DyadicRatCoreUp$ observation in $D$ to be read from those windows;
\item handoff and seal discipline, requiring the $\RegSeqRatUp$ rows $R$ and the $\RealUp$ seal row $S$ to be reached through $D$ and then transported by the displayed $H,C,P,N$ rows;
\item non-escape, requiring every public read to stay inside the displayed packet and to export no quotient stream, host real equality, countable-choice row, hidden limit selector, or ambient completeness theorem.
\end{enumerate}
\end{theorem}
\begin{proof}
Carrier exposure is projection from the tuple in \autoref{def:choice-free-limit-schedule-carrier}. The packet lists the precision request, schedule, windows, dyadic observations, regular-sequence handoff, real seal, transports, routes, provenance, and local naming row.

Request-to-schedule exactness is the first displayed route in $C$. Since the carrier has one schedule row $M$ for the request $E$, a consumer has no alternate coordinate from which to obtain a different schedule for the same precision request.

Window and dyadic readback coverage use the next two route segments. The route from $M$ to $W$ exposes only the finite window family, and the route from $W$ to $D$ makes the dyadic observation ledger depend on those windows rather than on an unbounded stream object.

Handoff and seal discipline are supplied by $R,S,H,C,P,N$. The regular-sequence handoff is reached after the dyadic ledger, the real seal is reached after the handoff, and transport or repacking must replay the same rows through $H,C,P,N$.

The non-escape clause follows from the carrier inventory. None of the displayed rows is a quotient stream, host real equality test, countable-choice row, hidden selector, or ambient completeness theorem, so the local naming certificate has no route that can export one.
\end{proof}

\begin{theorem}[Choice-free limit schedule non-escape]
\label{thm:choice-free-limit-schedule-nonescape}
Every $\RegSeqRatUp$, $\RealUp$, or completion-facing consumer of an accepted $\ChoiceFreeLimitScheduleUp$ packet factors through the finite route
$$
E,M,W,D,R,S,H,C,P,N.
$$
The consumer may inspect the precision request, schedule row, finite $\StreamNameUp$ windows, $\DyadicRatCoreUp$ observations, $\RegSeqRatUp$ handoff, $\RealUp$ seal, componentwise $\hsame$ transports, displayed $\Cont$ routes, $\Pkg$ provenance, and local $\NameCert$ row. It cannot inspect quotient streams, host equality of completed reals, countable choice, hidden limit selectors, or an ambient completeness theorem.
\end{theorem}
\begin{proof}
Start from the accepted carrier $L=(E,M,W,D,R,S,H,C,P,N)$. The only request-facing route begins at $E$ and reaches the schedule row $M$ through the displayed continuation coordinate.

Next read the finite observation surface. The schedule selects the windows in $W$, and the dyadic ledger $D$ is read from those windows. Thus a consumer cannot replace the displayed window family by a quotient stream or an unbounded stream-level selector.

Pass to the completion-facing rows. The handoff row $R$ and seal row $S$ are reached after the dyadic ledger and are repacked only through $H,C,P,N$, so the route into $\RegSeqRatUp$ and $\RealUp$ remains the same finite packet.

Finally apply the non-escape obligation of \autoref{thm:choice-free-limit-schedule-namecert-obligations}. The excluded analytic objects do not occur as coordinates of $L$, so no public read can export quotient streams, host equality of completed reals, countable choice, hidden limit selectors, or an ambient completeness theorem.
\end{proof}

\closureat{ChoiceFreeLimitScheduleUp}{obligationStr}

\begin{closurestatus}{\ChoiceFreeLimitScheduleUp}
  \constructivestory{A $\ChoiceFreeLimitScheduleUp$ packet is generated from a precision request, finite schedule row, StreamNameUp windows, DyadicRatCoreUp observations, RegSeqRatUp handoff, RealUp seal, componentwise transports, continuation routes, package provenance, and a local naming row.}
  \theoryclosure{\obligationClosure}
  \scopeclosed{\autoref{def:choice-free-limit-schedule-carrier}, \autoref{thm:choice-free-limit-schedule-namecert-obligations}, and \autoref{thm:choice-free-limit-schedule-nonescape} bind the finite precision-to-window-to-seal packet consumed by Real-completion routes.}
  \formalstatus{\unformalizedV}
  \bridgestatus{none}
  \origin{human}
  \notclaimed{This chapter does not close RealUp completeness, uniqueness of completed limits, quotient equality of streams, host real equality, a public standard bridge, countable choice, or Lean verification.}
  \upgradepath{Scoped closure requires binding the schedule classifier, transport stability, RegSeqRatUp handoff exactness, RealUp seal discipline, and consumer readback to the same displayed rows.}
\end{closurestatus}
